\documentclass[10pt,a4paper]{amsart}
\usepackage[margin=19mm]{geometry}
\usepackage{amsmath,amssymb,mathtools}
\usepackage[scr=rsfs,cal=euler]{mathalfa}
\usepackage{xfrac}
\usepackage[unicode,hidelinks,bookmarks=false]{hyperref}
\newcommand{\ie}{i.e.\ }
\newcommand{\cf}{cf.\ }
\newcommand{\resp}{respectively\ }
\newcommand{\Subsection}{\subsection}
\providecommand{\nobreakdash}{\nobreak\hbox{-}}
\setlength{\emergencystretch}{2em}
\multlinegap0pt
\usepackage{bm}
\usepackage{bbm}
\usepackage{dsfont}
\usepackage{xspace}
\usepackage{cancel}
\usepackage{mathtools}
\usepackage{amsthm}
\makeatletter
\@ifundefined{Lemma}{\newtheorem{Lemma}{Lemma}}{}
\makeatother
\title[Derivations of Lie Algebras of Vector Fields in Infinitely M]{Derivations of Lie Algebras of Vector Fields in Infinitely Many Variables}
\author{Oksana Bezushchak, Iryna Kashuba}
\date{Source: arXiv:2507.04541v1 (CC BY 4.0)}
\begin{document}
\maketitle

\noindent\emph{Source and licence.} Derivations of Lie Algebras of Vector Fields in Infinitely Many Variables. Version arXiv:2507.04541v1 (CC BY 4.0), distributed under the Creative Commons Attribution 4.0 International licence, as recorded in the arXiv licence metadata for this exact version and shown on the abstract page https://arxiv.org/abs/2507.04541v1. Full rights notice: licenses/arxiv-2507.04541-CC-BY-4.0.txt.\par\smallskip

\noindent\textit{Editorial scope.} Counted unit: the Lemma whose source label is \texttt{Lemma8} in the article (source0.tex lines 342--342, source label \texttt{Lemma8}), delivered together with the article's own complete original proof of it (source0.tex lines 344--488, one \texttt{proof} environment of 145 lines). No mathematics is written, repaired or re-derived: statement and proof are the article's own text line for line. Required context retained uncounted: Lemma6 (lines 277--277); Lemma7 (lines 289--293). Every definition, display and label of those lines is retained unchanged. Omitted from this package and not redistributed: 1--276, 278--288, 294--341, 343--343, 489--822. Nothing inside a retained excerpt is omitted or altered: every retained body is delivered exactly as the article's lines are written, so no omission receipt of any kind accompanies this package. Citations: the retained excerpts carry no citation command at all, so no bibliography entry of the article is transcribed and no reference is redistributed. Reference adaptation: none was needed and none was made; every reference inside the delivered unit resolves inside it, so no unresolved reference or citation is produced. Printed slips kept literal: no printed word, symbol, hypothesis or number of the counted statement or of its proof was altered. Layout and class: the article's own class is \texttt{amsart} with packages including amsmath, amssymb, hyperref, babel, cite; the wrapper is the lane's pinned \texttt{amsart} editorial wrapper at 10pt on A4, which restates the theorem-like environments the retained text needs and adds the article's own macro definitions that the retained text needs (none), with extra wrapper packages (bm, bbm, dsfont, xspace, cancel, mathtools). The delivered numbering is the unit's own. Assets: no retained span contains a figure environment, a graphics-inclusion command, a PDF-page inclusion, a table or a TikZ picture; no image file is copied into this package, the package declares no asset dependency and no image file is redistributed with it. Licence: version arXiv:2507.04541v1 of this article is distributed under the Creative Commons Attribution 4.0 International licence, as recorded in the arXiv licence metadata for this exact version and shown on the abstract page https://arxiv.org/abs/2507.04541v1. A full rights notice is delivered at licenses/arxiv-2507.04541-CC-BY-4.0.txt. Characters: every retained excerpt is pure ASCII, so the delivered engine, XeLaTeX with its default Latin Modern text font, reports no missing character.
	\begin{Lemma}\label{Lemma6} The first cohomology group vanishes: $$H^1\big(\mathfrak{sl}_{\infty}(\mathbb{F}), \widetilde{W}\big) = (0).$$ \end{Lemma}

	\begin{Lemma}\label{Lemma7} Let $w \in \widetilde{W}$, and $\mathfrak{sl}_\infty(\mathbb{F}) w \subseteq \widetilde{W}_X$. Then
		\[
		w = a \sum_{i=1}^{\infty} x_i \frac{\partial}{\partial x_i} + w',
		\]
		where $w' \in \widetilde{W}_X$, and for an arbitrary $i$, the partial derivative $\frac{\partial a}{\partial x_i}$ is a polynomial. \end{Lemma}

	\begin{Lemma}\label{Lemma8} The first cohomology group vanishes: $$H^1(L, \widetilde{W}_X) = (0).$$ \end{Lemma}

	\begin{proof} Let \( d : L \to \widetilde{W}_X \) be a derivation. Let's show that it is inner. By Lemma \ref{Lemma6}, there exists an element \( w \in \widetilde{W} \) such that
		\[
		d(x) = [w, x] \quad \text{for any } x \in \mathfrak{sl}_\infty(\mathbb{F}).
		\] By Lemma \ref{Lemma7}, $$ w = a \sum_{i=1}^\infty x_i \frac{\partial}{\partial x_i} + w',$$ where $ w' \in \widetilde{W}_X , $ and  $\partial a / {\partial x_i}$ is a polynomial for every  $i. $
		
		
		Take \( 1 \leq i \not= p < \infty \), $p\not= q.$ Then
		\[
		\left[ \frac{\partial}{\partial x_i}, x_p \frac{\partial}{\partial x_q} \right] = 0.
		\] Applying the derivation \( d \), we get:
		\[
		\left[ d\left( \frac{\partial}{\partial x_i}\right), \ x_p \frac{\partial}{\partial x_q} \right]
		+ \left[ \frac{\partial}{\partial x_i}, \Big[w, x_p \frac{\partial}{\partial x_q}\Big] \right]
		= \] \[\left[ d\left( \frac{\partial}{\partial x_i} \right) + \Big[\frac{\partial}{\partial x_i}, w\Big], \  x_p \frac{\partial}{\partial x_q} \right] = 0.
		\]
		
		Let
		\[
		d\left( \frac{\partial}{\partial x_i} \right) = \sum_{k=1}^\infty u_k \frac{\partial}{\partial x_k},\] where  $ u_k$  are polynomials.  We obtain: 
		\[
		\left[ \frac{\partial}{\partial x_i}, w \right] = \left[ \frac{\partial}{\partial x_i}, \ a\sum_{j=1}^\infty x_j \frac{\partial}{\partial x_j} +   w' \right]=
		\]
		\[ \left( \frac{\partial a}{\partial x_i}\right)  \sum_{j=1}^\infty x_j \frac{\partial}{\partial x_j} + a  \frac{\partial }{\partial x_i} + \left[ \frac{\partial}{\partial x_i}, w' \right].  \]
		
		
		Let
		\[
		\left[ \frac{\partial}{\partial x_i}, w' \right] = \sum_k v_k \frac{\partial}{\partial x_k}, \] where   $v_k$ are polynomials. The \( \frac{\partial}{\partial x_i} \)-component of the element \[   d\left(\frac{\partial}{\partial x_i}\right) +\left[  \frac{\partial}{\partial x_i}, w\right] \quad \text{is} \quad 
		\left(u_i + \left(\frac{\partial a}{\partial x_i}\right) x_i + a + v_i\right) \frac{\partial}{\partial x_i}.
		\]
		
		
		Suppose that the element $a$ is not a polynomial. Since \[  u_i + \left(\frac{\partial a}{\partial x_i} \right)x_i + v_i \] is a polynomial and every variable is involved in finitely many monomials in \( a \), it follows that there exists \( q \neq i \) such that
		\[
		\frac{\partial}{\partial x_q}\left(u_i + \left(\frac{\partial a}{\partial x_i}\right) x_i + a + v_i\right) \neq 0.
		\]
		
		Let $p \neq i,$ and $ p \neq q.$  Then the element
		\[
		\left[  d\left( \frac{\partial}{\partial x_i} \right) + \left[ \frac{\partial}{\partial x_i}, w \right] , \ x_p \frac{\partial}{\partial x_q}\right]
		\]
		has a nonzero \( \frac{\partial}{\partial x_i} \)-component. We got a contradiction. Thus, the element $a$ is a polynomial, and $w\in \widetilde{W}_X.$  
		
		
		Now, considering the derivation $d - \hat{w}$ instead of $d$, we assume that $d(\mathfrak{sl}_\infty(\mathbb{F})) = (0).$ Let
		\[
		d\left(\frac{\partial}{\partial x_i}\right) = \sum_{j=1}^{\infty} f_j \frac{\partial}{\partial x_j}, \] where each $ f_j $ is a polynomial in $X.$ Arguing as above, we see that for $p \neq i$ and $q \neq p$:
		\begin{equation}\label{eq5}
			\left[ \sum_{j=1}^{\infty} f_j  \frac{\partial}{\partial x_j}, \ x_p \frac{\partial}{\partial x_q} \right] = 
			f_p \frac{\partial}{\partial x_q} - x_p  \sum_{j=1}^{\infty} \ \frac{\partial f_j}{\partial x_q} \ \frac{\partial}{\partial x_j} = 0 .
		\end{equation}
		Fix $j \geq 1$. Choosing $q \neq j$, we get $\frac{\partial f_j}{\partial x_q} = 0.$  Hence,
		$f_j = f(x_j)$ is a polynomial in $x_j.$
		
		Now the right-hand side of \eqref{eq5} looks as
		\[
		f_p  \frac{\partial}{\partial x_q} - x_p  \frac{\partial f_p}{\partial x_q}\  \frac{\partial}{\partial x_p} = 0,
		\]
		which implies
		\[
		f_p(x_p) = x_p  \frac{\partial f_q(x_q)}{\partial x_q}.
		\] It implies that $f_q(x_q) = \alpha_q + \beta_q x_q$, for $\alpha_q, \beta_q \in \mathbb{F},$ and $f_p(x_p) =  \beta_q x_p.$  Hence, there exists $\beta \in \mathbb{F}$ such that for every $p\not= i, $  we have 
		$f_p = \beta x_p.$ Choosing $q = i$, we get $f_i = \alpha + \beta x_i$. Thus,
		\[
		d\left( \frac{\partial}{\partial x_i} \right) = \alpha_i  \frac{\partial}{\partial x_i} + \beta_i  \sum_{j=1}^{\infty} x_j  \frac{\partial}{\partial x_j}, \quad \text{where} \quad \alpha_i, \ \beta_i \in \mathbb{F}.
		\]
		
		Let $1 \leq i \neq j < \infty$. Applying the derivation $d$ to $[\frac{\partial}{\partial x_i}, \frac{\partial}{\partial x_j}] = 0$, we get:
		\[
		\left[ \alpha_i   \frac{\partial}{\partial x_i} + \beta_i  \sum_{k=1}^{\infty}  x_k  \frac{\partial}{\partial x_k}, \ \frac{\partial}{\partial x_j} \right] + \left[ \frac{\partial}{\partial x_i}, \ \alpha_j  \frac{\partial}{\partial x_j} + \beta_j  \sum_{k=1}^{\infty}  x_k  \frac{\partial}{\partial x_k} \right] = 
		\]
		\[ \beta_j  \frac{\partial}{\partial x_i} -\beta_i \frac{\partial}{\partial x_j} = 0,
		\]
		which implies $\beta_i = \beta_j = 0$.
		Hence, $$d\left(\frac{\partial}{\partial x_i}\right) = \alpha_i   \frac{\partial}{\partial x_i} \quad \text{for every} \quad  i \geq 1.$$
		
		
		Applying the derivation $d$ to
		\[
		\left[ \frac{\partial}{\partial x_i}, \ x_i  \frac{\partial}{\partial x_q} \right] = \frac{\partial}{\partial x_q}, \quad i \neq q,
		\] we get: $$ \alpha_i   \frac{\partial}{\partial x_q}= \alpha_q   \frac{\partial}{\partial x_q}.$$ Hence $ \alpha_i =\alpha$  for all $ i\geq 1. $ Now
		\[
		d\left( \frac{\partial}{\partial x_i} \right) = \alpha  \frac{\partial}{\partial x_i} = \left[ \frac{\partial}{\partial x_i}, \ \alpha \sum_{k=1}^{\infty}  x_k  \frac{\partial}{\partial x_k} \right].
		\]
		
		
		Replace \( d \) with $  d - \widehat{u} ,$ where $$ u= \alpha \sum_{k=1}^\infty x_k \frac{\partial}{\partial x_k}.$$
		Since
		\[
		\left[\mathfrak{sl}_\infty(\mathbb{F}), \ \sum_{k=1}^\infty x_k \frac{\partial}{\partial x_k}\right] = (0),
		\]
		we still have \( d(\mathfrak{sl}_\infty(\mathbb{F})) = (0) \),  and in addition to that,
		\[
		d\left( \frac{\partial}{\partial x_i} \right) = 0 \quad \text{for any } i \geq 1.
		\]
		
		We have:
		\[
		\mathfrak{gl}_\infty(\mathbb{F}) = \sum_{i=1}^\infty \mathbb{F} x_i \frac{\partial}{\partial x_i} + \mathfrak{sl}_\infty(\mathbb{F}).
		\] Let us prove that \[
		d\left( x_i \frac{\partial}{\partial x_i}\right) = 0.
		\] 
		
		For any \( p \geq 1 \), we compute
		\[
		\left[ \frac{\partial}{\partial x_p}, \ x_i \frac{\partial}{\partial x_i} \right] =
		\begin{cases}
			\frac{\partial}{\partial x_i}, & \text{if } p = i, \\
			0, & \text{if } p \neq i.
		\end{cases}
		\] In both cases,
		\[
		d\left( \left[ \frac{\partial}{\partial x_p}, \ x_i \frac{\partial}{\partial x_i} \right] \right) = 0.
		\]
		Since $$ d\left( \frac{\partial}{\partial x_p} \right) = 0,  \quad \text{then} \quad 
		\left[ \frac{\partial}{\partial x_p},\  d \left( x_i \frac{\partial}{\partial x_i} \right) \right] = 0.$$
		Assume
		\[
		d \left( x_i \frac{\partial}{\partial x_i} \right) = \sum_{k=1}^{\infty} a_k \frac{\partial}{\partial x_k}.
		\]
		For any \( k \geq 1 \), we have \( \frac{\partial}{\partial x_p} a_k = 0 \). Hence, \( a_k \in \mathbb{F} \). From the fact that 
		\[
		d \left( x_i \frac{\partial}{\partial x_i} - x_j \frac{\partial}{\partial x_j} \right) = 0,
		\]
		it follows that the element 
		\[
		\sum_{k = 1}^{\infty} a_k \frac{\partial}{\partial x_k}
		\]
		does not depend on the choice of the index  \( i \).
		
		Let \( q \ne i \). Then
		\[
		\left[ x_i \frac{\partial}{\partial x_i}, \ x_i \frac{\partial}{\partial x_q} \right] = x_i \frac{\partial}{\partial x_q}.
		\] Applying the derivation \( d \), we obtain:
		\[
		\left[ \sum_{k=1}^\infty a_k \frac{\partial}{\partial x_k}, \ x_i \frac{\partial}{\partial x_q} \right] = 0.
		\] Consequently, \( a_i = 0 \), which completes the proof that
		\[
		d(\mathfrak{gl}_\infty(\mathbb{F})) = (0).
		\]
		
		Thus, \[
		d(a) = \left[a, \ \alpha  \sum_{k=1}^\infty  x_k  \frac{\partial}{\partial x_k}\right] = 0 \quad \text{for all} \quad a \in \mathfrak{gl}_\infty(\mathbb{F}).
		\] The element $$\alpha  \sum_{k=1}^\infty x_k  \frac{\partial}{\partial x_k}$$ lies in $\widetilde{W}_{\text{fin}}.$ Thus, we proved that $d$ is an inner derivation. This completes the proof of the lemma. 
	\end{proof}

\end{document}
